\section{Methods}
\label{sec:methods}

This section fixes the training pipeline and the statistical
methodology.  The values reported here were committed to the
pre-registered plan in \texttt{paper/analysis\_plan.md} before any
main-result experiments were run; we note any deviations explicitly in
\Cref{sec:deviations}.

\subsection{Algorithm and Architecture}
\label{sec:algorithm}

We use RecurrentPPO from
\texttt{sb3-contrib}~\citep{sb3,schulman2017ppo,hausknecht2015drqn} for
all three treatments.  The architecture upstream of the LSTM differs
across treatments (\Cref{sec:obs-treatments}) but everything from the
LSTM onward is identical:
%
\begin{itemize}
\item \textbf{LSTM:} 1 layer, 256 hidden units (parameter-matched
  across all three treatments).
\item \textbf{Policy and value heads:} 2-layer MLPs of width 64 each,
  with \texttt{tanh} activations.
\item \textbf{Action distribution:} categorical over the 7 buttons.
\end{itemize}
%
The decision to vary only the encoder architecture is deliberate: any
performance difference between treatments is attributable to what the
observation stream makes available to the LSTM, not to differences in
recurrence or policy expressivity.

\subsection{Hyperparameters}
\label{sec:hyperparameters}

The RecurrentPPO hyperparameters used in the pilot are:
\todopilot{Confirm against the actual values in
\texttt{pokemon\_red\_ai/training/models.py} after the first run; the
values below are the pre-registered defaults.}
%
\begin{table}[h]
\centering
\small
\begin{tabular}{lr}
\toprule
Hyperparameter & Value \\
\midrule
Learning rate & $3 \times 10^{-4}$ \\
Batch size    & $256$ \\
Rollout length ($n\_\text{steps}$) & $2{,}048$ \\
Epochs per update & $10$ \\
Entropy coefficient & $0.01$ \\
Discount $\gamma$ & $0.99$ \\
GAE $\lambda$ & $0.95$ \\
Clip range & $0.2$ \\
Parallel envs & $1$ \\
\bottomrule
\end{tabular}
\end{table}
%
We deliberately use a single environment per training run rather than
the 8-env parallelism committed in the long-form plan, because the
PyBoy emulator is single-threaded per instance and parallelism on a
single machine yields diminishing returns past 4 envs in our
preliminary timing.  For the journal version we will revisit
parallelism once compute scales beyond a single workstation.

\subsection{Training Protocol}
\label{sec:protocol}

Each treatment is trained for $10^7$ environment steps across three
seeds (42, 123, 456).  No hyperparameter tuning is performed within
the pilot; all treatments use identical hyperparameters.  Training is
checkpoint-saved every 250{,}000 steps and the best checkpoint per run
(by training-time best-episode reward) is retained for evaluation.

\subsection{Evaluation Metric and Statistical Analysis}
\label{sec:stats}

The pre-registered \emph{primary} metric is Boulder Badge win-rate at
the end of training, evaluated as the fraction of 20 deterministic
rollouts that earn the badge before episode truncation.  At
10M~environment steps---one-tenth of the journal-version
budget---reaching Brock is a stretch goal rather than a baseline
expectation; in this paper we therefore report the primary metric
together with two informative secondary metrics:
%
\begin{itemize}
\item \textbf{Best episode reward} during training (accumulated event
  flags + discovery + time penalty).  This is a continuous proxy for
  progress and is what we use as the single number per seed when
  computing IQM.
\item \textbf{First episode at which each event flag is triggered.}
  Reported per-treatment as a milestone progression chart.
\end{itemize}

\paragraph{Aggregate metrics with confidence intervals.}
We use \texttt{rliable}~\citep{agarwal2021rliable} for all aggregate
statistics.  The primary point estimate per treatment is the
\emph{interquartile mean} (IQM) over the seed-level best-episode
rewards; we drop the bottom and top quartiles before averaging to
reduce sensitivity to outlier seeds.  Uncertainty is reported as a
95\% percentile bootstrap CI with $2{,}000$ resamples.

\paragraph{Probability of improvement.}
For pairwise comparisons we report the probability of improvement
(PoI), $\Pr[\text{score}_A > \text{score}_B]$, estimated via stratified
bootstrap.  We treat a treatment as superior to another only if the
95\% CI of PoI excludes 0.5.

\subsection{Deviations from the Pre-Registered Plan}
\label{sec:deviations}

The following deviations from the pre-registered plan are made
explicitly for the EWRL pilot and will be returned to nominal for the
journal version:
%
\begin{itemize}
\item \textbf{Seed count.}  Three seeds rather than the
  five-minimum committed for the main result.  We adopt three for the
  pilot to fit within the EWRL submission window; the journal version
  will use five (minimum) or eight (preferred) seeds per treatment.
\item \textbf{Training budget.}  $10^7$ steps rather than $10^8$.
  Reaching Brock is a stretch goal at this budget.  Effects detectable
  at $10^7$ are likely conservative estimates of effects at the
  full-budget version.
\item \textbf{Single environment per run.}  We run one parallel
  environment rather than the eight committed in the plan.  The
  journal version will revisit this once compute scales beyond a
  single workstation.
\end{itemize}
%
We do not deviate from the pre-registered hypotheses, primary or
secondary metrics, statistical tests, event-flag set, or evaluation
protocol.

\subsection{Compute and Reproducibility}
\label{sec:compute}

\todopilot{Fill from \texttt{paper/compute\_ledger.md} after pilots
finish: per-treatment wall-clock time, GPU utilisation peak, peak
memory, total GPU-hours.  Format as a small table.}

The training script, the evaluation harness, the statistical
pipeline, and the live training dashboards are open-sourced in the
project repository.  Training-run hyperparameter dumps and W\&B
artifact links for the runs reported here are provided in the
supplementary material.
